\section{Quantitative approach and reporting}
We describe each outcome on its own valid-response base. Separate core logit models relate the two outcomes to the three measures, age, grouped education and grouped recorded residence. The reference categories are completed secondary education and Bogotá. The residual residence group is neither a fourth city nor a sampling cluster. Extensions add grouped occupation (employee or worker as reference) and original first-arrival calendar year, not reconstructed tenure.

The five model rows cover the two core models, their two extensions and a separate refit of the remittance core model on the shared outcome-observed core base. The design checks that the shared base equals the domestic core base before omitting a duplicate domestic fit. We distinguish intended eligibility from the estimator's sample. Completeness and valid inputs do not establish random missingness or justify treating model-specific exclusions as deletions from the cohort.

Each core contrast compares average fitted probabilities at the higher and lower specified values of a measure, with other covariates at their observed values and the averages over that model's estimation sample. The endpoints are 0.50 and 0.75 for access/basic use and for the transfer ability/online-safety measure, and 0.50 and 1.00 for the two onboarding conditions. These are different increments. Observed support at the endpoints does not establish joint support for all combinations used in standardization.

The proposed figure reports differences in percentage points with pointwise 95\% delta-method intervals from a working robust covariance that assumes independent records. The six raw and Holm-adjusted p-values concern coefficient-null tests in the two core models, not simultaneous confidence intervals for the plotted probability contrasts. The working covariance does not account for dependence through unrecorded referral or household links. Cross-sectional associations cannot distinguish skills acquired through use from skills that precede it, and adjustment does not remove unmeasured confounding. This is a reporting specification, not a preregistration claim.

\subsection{Survey results: exhibits reserved}
The observed outcome distributions come first, then the two core models and the separately specified checks. The slots below contain no estimates, counts from outputs, or rendered figures. See the technical supplement for field, denominator and source definitions.

\exhibitplaceholder{T2/F1}{Reported disposition and domestic digital participation}{Each panel uses its own valid-response base. Neither panel is a share of money or a stage in a transaction funnel.}
\exhibitplaceholder{QT1}{Working model status and sample sizes}{Show intended and estimator sample sizes and reasons for unavailability by model. E.5 hold.}
\exhibitplaceholder{QF1}{Adjusted probability differences}{Six fixed core contrasts in percentage points with pointwise 95\% intervals. E.5 hold.}
\exhibitplaceholder{QT2}{Six core coefficient tests}{Report raw and Holm-adjusted Wald-test probabilities, not simultaneous intervals. E.5 hold.}
